% !TEX root = ../../main.tex
%% Corrigé du défi du chapitre 19 — vérifié par verifications/defis_16_20.py.

\begin{blocLivre}{cor-19-25a}
\begin{corrige}{25}{Le problème de Bâle : $\zeta(2)=\pi^{2}/6$}
\demarche{Les $m$ nombres $\cot^{2}\frac{r\pi}{2m+1}$ sont les racines d'un
polynôme explicite : Viète donne leur somme ; l'encadrement
$\cot^{2}\theta<\theta^{-2}<1+\cot^{2}\theta$ transmet cette somme exacte à
$\sum1/r^{2}$.}
\partieExo{A}{Une identité polynomiale}
\begin{questions}
  \item Moivre : $(\cos\theta+\mathrm{i}\sin\theta)^{2m+1}
    =\cos(2m+1)\theta+\mathrm{i}\sin(2m+1)\theta$. Par le binôme, la partie
    imaginaire provient des indices impairs $j=2k+1$, avec
    $\mathrm{i}^{2k+1}=(-1)^{k}\mathrm{i}$ : on obtient la formule.
  \item On divise par $\sin^{2m+1}\theta\neq0$ :
    $\frac{\cos^{2m-2k}\theta\,\sin^{2k+1}\theta}{\sin^{2m+1}\theta}
    =\bigl(\cot^{2}\theta\bigr)^{m-k}$, d'où
    $\res{\dfrac{\sin(2m+1)\theta}{\sin^{2m+1}\theta}=P_m\bigl(\cot^{2}\theta\bigr)}$.
  \item Pour $1\leqslant r\leqslant m$, $\theta_r=\frac{r\pi}{2m+1}\in\intOO{0}{\frac\pi2}$
    et $\sin\bigl((2m+1)\theta_r\bigr)=\sin r\pi=0$ : $P_m(x_r)=0$. La
    fonction $\cot^{2}$ est strictement décroissante sur $\intOO{0}{\frac\pi2}$ :
    les $x_r$ sont distincts. $P_m$ est de degré $m$, de coefficient
    dominant $\binom{2m+1}{1}=2m+1$ : ce sont toutes ses racines, et Viète
    donne
    \[
    \sum_{r=1}^{m}x_r=\frac{\binom{2m+1}{3}}{2m+1}=\frac{(2m)(2m-1)}{6}
    =\res{\frac{m(2m-1)}{3}} .
    \]
\end{questions}
\partieExo{B}{L'encadrement}
\begin{questions}[resume]
  \item $\theta-\sin\theta$ et $\tan\theta-\theta$ sont nulles en $0$ et de
    dérivées $1-\cos\theta>0$ et $\tan^{2}\theta>0$ sur
    $\intOO{0}{\frac\pi2}$ : $0<\sin\theta<\theta<\tan\theta$. En passant à
    l'inverse des carrés :
    $\cot^{2}\theta<\frac{1}{\theta^{2}}<\frac{1}{\sin^{2}\theta}=1+\cot^{2}\theta$.
  \item Avec $\frac1{\theta_r^{2}}=\frac{(2m+1)^{2}}{\pi^{2}r^{2}}$ et en
    sommant :
    $\frac{m(2m-1)}{3}<\frac{(2m+1)^{2}}{\pi^{2}}\sum_{r=1}^{m}\frac1{r^{2}}
    <m+\frac{m(2m-1)}{3}=\frac{2m(m+1)}{3}$.
  \item En multipliant par $\frac{\pi^{2}}{(2m+1)^{2}}$, les deux bornes
    $\frac{\pi^{2}m(2m-1)}{3(2m+1)^{2}}$ et
    $\frac{2\pi^{2}m(m+1)}{3(2m+1)^{2}}$ tendent vers
    $\frac{2\pi^{2}}{12}$ : par encadrement,
    $\res{\zeta(2)=\dfrac{\pi^{2}}{6}}$.
\end{questions}
\end{corrige}
\end{blocLivre}

\begin{blocLivre}{cor-19-25b}
\begin{corrigeSuite}
\partieExo{C}{Conséquences}
\begin{questions}[start=7]
  \item Les termes pairs donnent
    $\sum\frac{1}{(2k)^{2}}=\frac{\zeta(2)}{4}$, donc les impairs
    $\frac34\zeta(2)=\res{\frac{\pi^{2}}{8}}$ ; et la série alternée vaut
    $\frac{\pi^{2}}{8}-\frac{\pi^{2}}{24}=\res{\frac{\pi^{2}}{12}}$
    (sommation par paquets de séries absolument convergentes).
  \item $\zeta(2)\in\intOO{S_{10}+\frac1{11}}{S_{10}+\frac{1}{10{,}5}}
    =\intOO{1{,}640\,677}{1{,}645\,006}$, et $\pi=\sqrt{6\zeta(2)}$ :
    \[
    \res{3{,}1375<\pi<3{,}1417},
    \]
    encadrement d'amplitude $0{,}0041$ avec dix termes seulement — grâce au
    contrôle fin du reste. (La vraie valeur $3{,}14159\dots$ est bien dans
    l'intervalle.)
\end{questions}
\end{corrigeSuite}
\end{blocLivre}
